% Generated draft: requires mathematical and source review.
\documentclass[11pt]{article}
\usepackage[margin=25mm]{geometry}
\usepackage{fontspec}
\setmainfont{Latin Modern Roman}
\usepackage{amsmath,amssymb,mathtools,unicode-math}
\setmathfont{Latin Modern Math}
\usepackage{enumitem,needspace,xurl,hyperref}
\setlength{\parindent}{0pt}
\setlength{\parskip}{5pt}
\setlength{\emergencystretch}{3em}
\newcommand{\N}{\mathbb{N}}
\newcommand{\Z}{\mathbb{Z}}
\newcommand{\Q}{\mathbb{Q}}
\newcommand{\R}{\mathbb{R}}
\newcommand{\C}{\mathbb{C}}
\newcommand{\F}{\mathbb{F}}
\newcommand{\A}{\mathbb{A}}
\newcommand{\PP}{\mathbb{P}}
\newcommand{\E}{\mathbb{E}}
\newcommand{\eps}{\varepsilon}
\newcommand{\dd}{\,\mathrm{d}}
\begin{document}
\section*{UnsolvedMath ID: 30002107}
\textbf{Source problem code:} catalogue entry
\subsection{Definitions and mathematical statement}
Let \(K\) denote the complete discretely valued base field mentioned in the record, and let \(G\) be a connected reductive group over that field. Write \(\mathcal O_K\) for its valuation ring. For terminology only, if \(E\) is a finite-dimensional \(K\)-vector space, an \(\mathcal O_K\)-lattice in \(E\) is a finitely generated \(\mathcal O_K\)-submodule \(L\subseteq E\) such that \(K L=E\). A lattice chain is a nonempty family of such lattices totally ordered by inclusion. These definitions do not prescribe indexing, periodicity, or additional compatibility conditions.

The supplied source asks the following construction question: given \(G\) and a representation \(V\), which lattice chains in the space written \(V\otimes K\) should be used to construct a Tannakian description of the Bruhat--Tits building of \(G\) and its parahoric group schemes?

This remains an incompletely specified construction request. The record does not identify the field or ring over which \(V\) is initially defined, the base of the displayed tensor product, the dimension or faithfulness requirements on \(V\), or the required class of lattice chains. It also does not specify which version of the building is intended, which parahoric group schemes are to be recovered, or the precise correspondence or equivalence constituting a Tannakian description. In particular, it is unclear whether the intended data concern one representation or a tensor-compatible assignment across representations. Consequently, the supplied evidence does not determine a precise existence assertion or characterization theorem.
\subsection{Short English statement}
The question asks which lattice chains attached to representations of a connected reductive group should describe its Bruhat--Tits building and parahoric group schemes through Tannakian methods. Essential input conventions and the required output correspondence are unspecified, so this draft preserves the construction question and records the missing information.
\subsection{Sources}
\begin{itemize}
\item Supplied local UnsolvedMath registry record, permanent ID \(30002107\), title: Tannakian Lattice Chains for Bruhat--Tits Buildings; exact source problem code: \texttt{catalogue entry}. Collection: Oberwolfach Reports Open Problems. The registry supplies proposed year 2012 and names no proposer.
\item Supplied report citation: \emph{Reductions of Shimura Varieties (2013)}, Oberwolfach Reports, DOI \texttt{10.4171/owr/2012/32}, pp. 1961--2011. Source URL and supplied literature source labelled ``original OWR report'': \url{https://doi.org/10.4171/owr/2012/32}. The background describes these reports as documenting mathematical workshops and their open-problem discussions.
\item Literature assessment dated exactly \texttt{2026-08-21}: the supplied curation labels the record open, but says that the extracted sentence merely requests a choice of lattice chains and does not specify a mathematical assertion; the underlying general parahoric construction is therefore unresolved or unclear from this record. This is supplied dated curation, not an independently verified assessment of the literature or current mathematical status.
\item The registry labels the statement \texttt{corrected\_verified} and reports that its mathematical content was checked against the official Oberwolfach report, with notation, context, and question form normalized without changing the source-level claim. This verification claim is supplied metadata and has not been independently checked here.
\end{itemize}
\end{document}
